\clearpage\thispagestyle{empty}
\pdfbookmark[0]{Sources for the review of Klüber on the Mannheim Observatory}{klueber1811-sources}
\begin{center}\textbf{Sources for the review of Klüber on the Mannheim Observatory}\end{center}
\noindent This AI-translated review follows Carl Friedrich Gauss, \textit{Werke}, vol.~VI (Göttingen, 1874), pp.~536--537. The review first appeared under Mannheim and Heidelberg in \textit{Göttingische gelehrte Anzeigen}, nos.~192 and 193, 2 December 1811, pp.~1913--1916.
\par\smallskip\noindent The collected edition prints \textit{Astonomen} for \textit{Astronomen} and omits the full stop before the sentence about the bombardment. It uses superscript m and s in the longitude from Paris, where the first printing uses primes.
\par\smallskip\noindent Johann Ludwig Klüber's \href{https://doi.org/10.11588/diglit.19120}{\textit{Die Sternwarte zu Mannheim}} (1811) describes a temporary \textit{Mire méridienne}, a target for aligning a transit telescope, in the passage at \href{https://doi.org/10.11588/diglit.19120#0061}{scan~0061} of the Heidelberg University Library copy.
\par\smallskip\noindent \href{https://gdz.sub.uni-goettingen.de/id/PPN236007467}{Collected text}. First printing, \href{https://images.p.sub.uni-goettingen.de/iiif/image/gdz:PPN319721507_1811:00001947/full/1650,/0/default.jpg}{p.~1913}, \href{https://images.p.sub.uni-goettingen.de/iiif/image/gdz:PPN319721507_1811:00001948/full/1650,/0/default.jpg}{p.~1914}, \href{https://images.p.sub.uni-goettingen.de/iiif/image/gdz:PPN319721507_1811:00001949/full/1650,/0/default.jpg}{p.~1915} and \href{https://images.p.sub.uni-goettingen.de/iiif/image/gdz:PPN319721507_1811:00001950/full/1650,/0/default.jpg}{p.~1916}.
\par\smallskip\noindent Translation and typesetting produced by AI.
\clearpage
